% =====================================================================
% Слайд 12 — Предлагаемая архитектура (РИС. 1 тезисов)
% =====================================================================
\begin{frame}{Предлагаемая архитектура (РИС. 1 тезисов)}
  \footnotesize

  \rlhi{Идея:} не заменять классический трекер, а подвесить
  \emph{рядом} лёгкий RL-агент. Базовый трекер остаётся ядром.
  Агент читает состояние трекера и историю темпа, выдаёт темповый
  коэффициент $a_t \in [0.5, 2.0]$ на горизонте 200--500\,мс.

  \begin{columns}[T,onlytextwidth]
    \column{0.50\textwidth}
    \centering
    \resizebox{0.92\linewidth}{!}{%
      % СХЕМА 8 — Предлагаемая архитектура (РИС. 1 тезисов)
\begin{tikzpicture}[
  every node/.style={font=\footnotesize},
  box/.style={draw, thick, minimum height=0.95cm, minimum width=2.0cm,
              align=center, rounded corners=1.5pt, fill=white, inner sep=3pt},
  rlbox/.style={box, line width=1.0pt, fill=black!6},
  arr/.style={-{Latex[length=2mm]}, thick},
]
  \node[box] (tr) at (0, 0)
       {\textbf{Базовый трекер}\\\scriptsize HSMM + OLTW + anchor};

  \node[rlbox] (pi) at (0, 2.1)
       {\textbf{Политика} $\pi_\theta$\\\scriptsize MLP};
  \draw[arr] (tr) -- node[right, font=\scriptsize] {$s_t$} (pi);

  \node[box] (env) at (3.6, 0)
       {\textbf{Среда}\\\scriptsize (рендерер)};
  \draw[arr] (pi.east) -| node[above, near start, font=\scriptsize]
       {$a_t$ (темп)} (env.north);

  \node[rlbox] (rew) at (3.6, -2.1)
       {\textbf{Награда} $r_t$\\\scriptsize формула (1)};
  \draw[arr] (env) -- (rew);
  \draw[arr] (rew.west) -| (pi.south);
\end{tikzpicture}
%
    }

    \column{0.50\textwidth}
    \textbf{Новая ниша:}
    \begin{itemize}\setlength{\itemsep}{0.35em}
      \item \emph{Dorfer 2018} --- RL заменяет трекер целиком.
      \item \emph{Peter 2023} --- RL для symbolic alignment.
      \item \emph{RL-Duet, ReaLchords} --- RL \emph{генерирует} ноты.
      \item \emph{Наш подход} --- RL поверх работающего трекера,
            \emph{только} темп, без замены классики.
    \end{itemize}

    \vspace{0.4em}
    {\itshape\color{black!75}
     Если RL-агент молчит --- система деградирует до проверенного
     гибрида, а не ломается.}
  \end{columns}
\end{frame}
